% 
%
% (c) 2025 Autor, ETH Zürich
%
% !TEX root = linalg_I-II_summary.tex
% !TEX encoding = UTF-8
%

\section{Bilinear \& Quadratic Forms}
\begin{definition}{Bilinear Form}{bilin_form}
	Let $V$ be a vector space over $K$. A function $B: V \times V \to K$ is called a \textbf{bilinear form} if the following two conditions are satisfied:
	\begin{enumerate}
		\item $\forall \, w \in V$ the map $B(\cdot , w):V \to K, \quad v \mapsto B(v, w)$ is linear.
		\item $\forall \, v \in V$ the map $B(v, \cdot):V \to K, \quad w \mapsto B(v, w)$ is linear.
	\end{enumerate}
\end{definition}

\begin{lemma}{Matrix Representation of a Bilinear Form}{mat_bilin_form}
	Let $V$ be a vector space over $K$ and $\mathcal{C} = (e_1, \hdots , e_n)$ a basis for $V$. Let $B: V \times V \to K$ be a bilinear form. Then $\exists \, [B]_{\mathcal{C}} \in M_{n \times n}(K)$ s.t.
	\[
		B\Big(\sum_{i=1}^{n} c_ie_i, \sum_{j=1}^{n}d_je_j\Big) = 
		\begin{pmatrix}
			c_1 & \hdots & c_n
		\end{pmatrix}
		\cdot [B]_{\mathcal{C}} \cdot
		\begin{pmatrix}
			d_1\\
			\vdots\\
			d_n
		\end{pmatrix} \qquad \forall \, c_i, d_j \in K.
	\]
	The matrix $[B]_{\mathcal{C}}$ is called the matrix representation of $B$ in the basis $\mathcal{C}$.
\end{lemma}

We saw also that for $T\in \operatorname{End}(V)$ and $\mathcal{C}$ a basis for $V$ we have a matrix representation $[T]^{\mathcal{C}}_{\mathcal{C}} \in M_{n\times n}(K)$, and that if $\mathcal{D}$ is another basis for $V$, then
\begin{equation}
	\label{eq:trans_mat}
	[T]^{\mathcal{D}}_{\mathcal{D}} = \underbrace{[id_V]^{\mathcal{C}}_{\mathcal{D}}}_{\displaystyle = \big([id_V]^{\mathcal{D}}_{\mathcal{C}}\big)^{-1}} \cdot [T]^{\mathcal{C}}_{\mathcal{C}} \cdot [id_V]^{\mathcal{D}}_{\mathcal{C}}.
\end{equation}
(Corollary \ref{cor*change_base_endo}).
What happens for bilinear forms?

\begin{lemma}{Change of Bases for a Bilinear Form}{change_base_bilin}
	Let $V$ be a finite dimensional vector space over $K$ and $\mathcal{C}, \mathcal{D}$ be two bases for $V$. Let $B:V \times V \to K$ be a bilinear form. Then
	\[
		[B]_{\mathcal{D}} = \big([id_V]^{\mathcal{D}}_{\mathcal{C}}\big)^T\cdot [B]_{\mathcal{C}} \cdot [id_V]^{\mathcal{D}}_{\mathcal{C}}.
	\]
	(This is quite different from Equation \eqref{eq:trans_mat}). 
\end{lemma}

\begin{definition}{Symmetric/Positive Definite Bilinear Form \& Quadratic Form}{symm_bilin_form}
	\begin{enumerate}
		\item A bilinear form $B:V \times V\to K$ is called \textbf{symmetric} if
		\[
			\forall \, v,w \in V : \; B(v, w) = B(w, v).
		\]
		\item Assume that $K = \mathbb{R}$. A bilinear form $B:V \times V \to K$ is called \textbf{positive definite} if 
		\[
			\forall \, v\in V:\; B(v, v) > 0.
		\]
		\item For a bilinear form $B:V \times V \to K$, we define
		\[
			q_B:V \to K, \quad q_B(v) := B(v, v). 
		\]
		We call $q_B$ the \textbf{quadratic form} associated to $B$. Sometimes we just denote it by $q$.
	\end{enumerate}
\end{definition}

\begin{remark}
	Note that $B \in M_{n\times n}(K)$ and $B^T \in M_{n\times n}(K)$ give the same quadratic form $q:K^n \to K$. Indeed, $\forall \, v\in K^n$
	\[
		q_B(v) = v^T B \, v = (v^T B \, v)^T = v^T B^T v = q_{B^T}(v),
	\]
	where we used that $q_B(v)$ is a scalar.
\end{remark}

Assume now that $2 \neq 0$ in $K$, then (i.e., $1 + 1 \neq 0$). Then also $\tfrac{1}{2}(B + B^T)$ gives the same quadratic form as $B$.

\textbf{Conclusion:} if $2 \neq0$ in $K$, then for the discussion on quadratic forms $q:K^n \to K$, we can assume they all come from symmetric matrices / bilinear forms.

\begin{proposition}{Polarization Identity}{pol_id}
	Assume $2 \neq 0$ in $K$. (note that then also $4 \neq 0$.) If $B$ is symmetric and defines $q := q_B$, then $\forall \, v, w \in V$ we have
	\begin{enumerate}
		\item 
		\[
			B(v, w) = \frac{1}{4} q(v + w) - \frac{1}{4} q(v - w).
		\]
		\item 
		\[
			B(v, w) = \frac{1}{2}\big(q(v + w) - q(v) - q(w)\big).	
		\]
	\end{enumerate}
\end{proposition}






